                                         8th IMC 2001
                                        July 19 - July 25
                                     Prague, Czech Republic

   First day

    Problem 1.
    Let n be a positive integer. Consider an n×n matrix with entries 1, 2, . . . , n2
written in order starting top left and moving along each row in turn left–to–
right. We choose n entries of the matrix such that exactly one entry is chosen
in each row and each column. What are the possible values of the sum of the
selected entries?

   Solution. Since there are exactly n rows and n columns, the choice is of
the form
                         {(j, σ(j)) : j = 1, . . . , n}
where σ ∈ Sn is a permutation. Thus the corresponding sum is equal to
                   n
                   X                             n
                                                 X              n
                                                                X          n
                                                                           X
                         n(j − 1) + σ(j) =               nj −         n+         σ(j)
                   j=1                           j=1            j=1        j=1

             n           n           n
             X           X           X                    n(n + 1)        n(n2 + 1)
        =n         j−          n+          j = (n + 1)             − n2 =           ,
             j=1         j=1         j=1
                                                             2                2

   which shows that the sum is independent of σ.

   Problem 2.
   Let r, s, t be positive integers which are pairwise relatively prime. If a and b
are elements of a commutative multiplicative group with unity element e, and
                t
ar = bs = (ab) = e, prove that a = b = e.
   Does the same conclusion hold if a and b are elements of an arbitrary non-
commutative group?

   Solution. 1. There exist integers u and v such that us + vt = 1. Since
ab = ba, we obtain
                                 v
          us+vt       us        t          us       us          u
ab = (ab)       = (ab)     (ab)      = (ab) e = (ab) = aus (bs ) = aus e = aus .
                                             r                    us
Therefore, br = ebr = ar br = (ab) = ausr = (ar )                      = e. Since xr + ys = 1 for
suitable integers x and y,

                                b = bxr+ys = (br )x (bs )y = e.

It follows similarly that a = e as well.
    2. This is not true. Let a = (123) and b = (34567) be cycles of the permu-
                                                                     7
tation group S7 of order 7. Then ab = (1234567) and a3 = b5 = (ab) = e.
                                                 ∞
                                                 X  tn
   Problem 3. Find lim (1 − t)                           , where t % 1 means that t ap-
                               t%1
                                              n=1
                                                  1 + tn
proaches 1 from below.

                                                     1
     Solution.
                              ∞                                    ∞
                              X    tn           1−t               X     tn
               lim (1 − t)            n
                                        =  lim         · (− ln t)            =
           t→1−0
                              n=1
                                  1+t     t→1−0 − ln t
                                                                  n=1
                                                                      1 + tn


                       ∞                           ∞             Z ∞
                       X        1                 X      1             dx
     = lim (− ln t)              −n ln t
                                         =  lim h           nh
                                                               =            = ln 2.
       t→1−0
                       n=1
                           1 + e           h→+0
                                                  n=1
                                                      1 + e       0  1 + ex


   Problem 4.
   Let k be a positive integer. Let p(x) be a polynomial of degree n each of
whose coefficients is −1, 1 or 0, and which is divisible by (x − 1)k . Let q be a
prime such that lnq q < ln(n+1)
                           k
                                . Prove that the complex qth roots of unity are
roots of the polynomial p(x).

     Solution. Let p(x) = (x−1)k ·r(x) and εj = e2πi·j/q (j = 1, 2, . . . , q −1). As
is well-known, the polynomial xq−1 + xq−2 + . . . + x + 1 = (x − ε1 ) . . . (x − εq−1 )
is irreducible, thus all ε1 , . . . , εq−1 are roots of r(x), or none of them.
                                                                         Qq−1
     Suppose that none of ε1 , . . . , εq−1 is a root of r(x). Then j=1 r(εj ) is a
rational integer, which is not 0 and
                                   q−1
                                   Y                   q−1
                                                       Y                     q−1
                                                                             Y
                (n + 1)q−1 ≥                p(εj ) =         (1 − εj )k ·          r(εj ) ≥
                                   j=1                 j=1                   j=1

                                       k
                     q−1
                     Y
                ≥          (1 − εj )       = (1q−1 + 1q−2 + . . . + 11 + 1)k = q k .
                     j=1

This contradicts the condition lnq q < ln(n+1)
                                          k
                                               .

   Problem 5.
   Let A be an n × n complex matrix such that A 6= λI for all λ ∈ C. Prove
that A is similar to a matrix having at most one non-zero entry on the main
diagonal.

    Solution. The statement will be proved by induction
                                                          onn. For n = 1,
                                                        a b
there is nothing to do. In the case n = 2, write A =          . If b 6= 0, and
                                                        c d
c 6= 0 or b = c = 0 then A is similar to
                                                           
               1 0       a b          1  0            0      b
                                             =
              a/b 1      c d        −a/b 1        c − ad/b a + d
or                                                                                
                    1 −a/c             a     b        1 a/c                0 b − ad/c
                                                                   =                         ,
                    0  1               c     d        0 1                  c   a+d
respectively. If b = c = 0 and a 6= d, then A is similar to
                                                       
                   1 1       a 0        1 −1          a d−a
                                                 =              ,
                   0 1       0 d        0 1           0     d

                                                       2
and we can perform the step seen in the case b 6= 0 again.
                                                                          0
   Assume
      0 now      that n > 3 and the problem has been solved for all n < n. Let
       A ∗
A=                 , where A0 is (n − 1) × (n − 1) matrix. Clearly we may assume
        ∗ β n
                                                                           
                                                                        0 ∗
that A0 6= λ0 I, so the induction provides a P with, say, P −1 A0 P =           .
                                                                        ∗ α n−1
But then the matrix
                   −1        0                  −1 0             
                    P      0     A ∗         P 0          P AP ∗
           B=                                        =
                      0    1      ∗ β        0 1               ∗      β
is similar to A and   its diagonal
                                     is (0, 0, . . . , 0, α, β). On the other hand, we may
                       0 ∗
also view B as                   , where C is an (n − 1) × (n − 1) matrix with diagonal
                       ∗ C n
(0, . . . , 0, α, β). If the inductive
                                       hypothesis is applicable to C, we would have
  −1                              0 ∗
Q CQ = D, with D =                             so that finally the matrix
                                  ∗ γ n−1
                                                                                            
          1  0                  1 0                1  0             0 ∗        1 0               0 ∗
E=                    ·B·                 =                                              =
          0 Q−1                 0 Q                0 Q−1            ∗ C        0 Q               ∗ D

  is similar to A and its diagonal is (0, 0, . . . , 0, γ), as required.
  The inductive argument can fail only when n − 1 = 2 and the resulting
matrix applying P has the form
                                                        
                                          0 a b
                           P −1 AP =  c d 0 
                                          e 0 d

where d 6= 0. The numbers a, b, c, e cannot be 0 at the same time. If, say,
b 6= 0, A is similar to
                                                             
        1 0 0         0 a b     1 0 0               −b     a    b
     0 1 0  c d 0  0 1 0  =                   c     d    0 .
        1 0 1         e 0 d    −1 0 1            e−b−d a b+d

Performing half of the induction step again, the diagonal of the resulting matrix
will be (0, d − b, d + b) (the trace is the same) and the induction step can be
finished. The cases a 6= 0, c 6= 0 and e 6= 0 are similar.

   Problem 6.
   Suppose that the differentiable functions a, b, f, g : R → R satisfy

               f (x) ≥ 0, f 0 (x) ≥ 0, g(x) > 0, g 0 (x) > 0 for all x ∈ R,

      lim a(x) = A > 0,           lim b(x) = B > 0,                 lim f (x) = lim g(x) = ∞,
      x→∞                        x→∞                             x→∞           x→∞

and
                                  f 0 (x)        f (x)
                                          + a(x)       = b(x).
                                  g 0 (x)        g(x)
Prove that
                                                  f (x)        B
                                      lim                 =       .
                                      x→∞ g(x)                A+1

                                                      3
    Solution. Let 0 < ε < A be an arbitrary real number. If x is sufficiently
large then f (x) > 0, g(x) > 0, |a(x) − A| < ε, |b(x) − B| < ε and


                             f 0 (x)        f (x)   f 0 (x)           f (x)
(1)        B − ε < b(x) =            + a(x)       <         + (A + ε)       <
                             g 0 (x)        g(x)    g 0 (x)           g(x)
                                      A                   A−1 0
           (A + ε)(A + 1) f 0 (x) g(x) + A · f (x) · g(x)      · g (x)
         <               ·                         A                  =
                 A                   (A + 1) · g(x) · g 0 (x)
                                                     A 0
                       (A + ε)(A + 1)   f (x) · g(x)
                     =                ·         A+1 0 ,
                             A              g(x)

thus
                                       A 0
                          f (x) · g(x)         A(B − ε)
(2)                               A+1 0 > (A + ε)(A + 1) .
                              g(x)

    It can be similarly obtained that, for sufficiently large x,
                                      A 0
                         f (x) · g(x)           A(B + ε)
(3)                              A+1 0 < (A − ε)(A + 1) .
                             g(x)

      From ε → 0, we have
                                                  A 0
                                     f (x) · g(x)        B
                           lim               A+1 0 = A + 1 .
                          x→∞
                                         g(x)

By l’Hospital’s rule this implies
                                                 A
                        f (x)       f (x) · g(x)      B
                    lim       = lim          A+1 =      .
                   x→∞ g(x)    x→∞                  A +1
                                        g(x)




                                              4
